\subsection{Results for Heterogeneity Analysis}
\label{subsec:results_heterogeneity}

\begin{customltable}
    \caption{Heterogeneity Results for different Subgroups of Settlement Schemes}
    \label{tab:base_heterogeneity_results}
    \setlength{\tabcolsep}{10pt}
    \small
    \tablerendering{tableBase_Heterogeneity}

    \tabnotes{This table shows a subsample of the Baseline Spatial RDD estimates at a 3000m bandwidth, including covariates and unknown settlements.
    Treatment effects, std. errors and means are multiplied by 100 for better readability.
    Columns (1) and (2) split the sample into conventional and squatter schemes, based on classification by \cite{boonePromisedLandSettlement2021}.
    Columns (3) and (4) split the sample into low and high population density settlements, based on the median population density per hectare of all settlements.
    Columns (5) and (6) split the sample into fast and slow registration of allocated settlements, based on the median time between allocation and registration of all settlements.
    Column (7) shows the specific results for the agricultural settlement schemes with a long registration period by restricting the sample to conventional, low density and slow registration settlements.
    Point estimates are conventional, with robust standard errors clustered at the 1200m grid level. *** $p < 0.01$, ** $p < 0.05$, * $p < 0.1$.}
\end{customltable}


\cref{tab:base_heterogeneity_results} uses the same data as the main results table, shown in the same two-panel design.
The focus of this table lies in using subsets of the main data to further strengthen the explanatory power of the model.
Significance numbers as high as in \cref{tab:base_baseline_results} could not be achieved, but significant effects are consistent with observations from the previous chapter.
All regressions have similar model specifications, with each using the MSE-optimal bandwidth,
and the covariates elevation, slope, distance to the nearest colonial railroad and distance to the nearest colonial city.
Additionally, the settlements with unknown boundaries have been excluded from all specifications.

\paragraph{Scheme Type.}
Columns (1) and (2) differentiate between conventional and squatter settlements.
The assignment to a conventional scheme reduces both indices significantly, the NDVI by $0.00769$ and the NDMI by $0.00947$.
When assigned to a squatter scheme, no significant change can be reported.
This is expected, as squatter schemes are often allocated in less fertile and forested areas.
However, the great difference in the number of observations could also lead to this outcome.
This difference is due to the fact that almost none of the schemes allocated in the region of the Kenya Highlands are squatter schemes.
With fewer observations, the standard error increases which in turn could reduce significance.

\paragraph{Population Density.}
Columns (3) and (4) differentiate between low density and high density settlements.
This can be used to further test the explanatory power of the first two columns as most conventional schemes had a low population density suitable for agricultural land use.
Analogous to column (1), the assignment to a low density scheme reduces the NDVI significantly by $0.00868$ and the NDMI by $0.01075$.
In contrast to insignificant squatter schemes, the allocation of land to a high density scheme causes a reduction of $0.00681$ in NDMI, significant at the 10\% level.
While this is an interesting change, this could be due to the High/Low cutoff using a median split, which in turn did balance the observations.

\paragraph{Registration Period.}
The next two columns (5) and (6) compare fast registration periods and slow registration periods.
This is used to test for Hypothesis 2.
First, the fast registration period has a negative impact of $0.00792$ on NDMI, once again at the 10\% level.
Interestingly, the schemes with slow registration periods experience significant NDVI and NDMI reductions at the 5\% level by $0.00719$ and $0.00955$, with standard errors reduced by about 10\%.
Observations split almost evenly between the two subsets.
To test for Hypothesis 2, it is necessary to conduct a $Z$-test to compare the two coefficients.
$$Z = \frac{TE_6 - TE_5}{\sqrt{SE^2_6 + SE^2_5}} \approx -0.29$$
This $Z$-score is not significant, and while the negative impact of settlement is significant, the null hypothesis of equal coefficients cannot be rejected.
As a result, it is not possible to draw conclusions about the differences between the two registration periods, and Hypothesis 2 cannot be confirmed.

\begin{customfigure}
    \caption{Regression Discontinuity Plot for Column (7) of \cref{tab:base_heterogeneity_results}}
    \includegraphics[width=0.6\textwidth]{figureBase_NDMI_HCOMBI}
    \label{fig:base_ndmi_hcombi}
    
    \fignotes{This figure shows the LATE on NDMI aggregated for the time period 2023 to 2025.
    Observations used for the regression are a subset of the main dataset, including only conventional, low density and slow registration settlements.
    Dots represent data aggregated in 300m bins for visual representation, the regression line is based on the full dataset.
    The slightly visible gray shaded areas represent the 90\% confidence intervals.}
\end{customfigure}

\paragraph{Combined subset of conventional, low density and slow registration.}
Column (7) shows the results for a combined subset of conventional, low density and slow registration settlements.
For a better understanding of how this looks from a geographical perspective,
\cref{fig:base_ndmi_hcombi} shows the development of NDMI in both the treatment and control groups, with the cutoff line at 0m.
The assignment to this subset reduces the NDVI by $0.00864$ and the NDMI by $0.01206$, significant at the 10\% and 5\% level respectively.
This is the highest reduction in NDMI of all subsets, and the second-highest reduction in NDVI.
However, the MSE-optimal bandwidth for this subset is larger compared to the other subsets, 
which could introduce effects into the model that are not present at the cutoff line and in other subsets.

While once again, many of the results are significant, the treatment effects still remain small in magnitude.
This will be further discussed in \cref{subsec:limitations}.